% GENERATED FILE. Edit canonical lessons, then run npm run generate:books.
% canonical-source-hash: fnv1a64:11bbb725abde7a12
% canonical-lessons: SA-C328-vratam, SA-C328-upavasah, SA-C328-naivedyam, SA-C328-arpayati, SA-C328-prasadah

\chapter{Fasting and Offerings}
\label{ch:sa-fasting-and-offerings}
\hlchaptermodality{328}

\begin{chapteropening}
\textbf{By the end of this chapter:} \emph{I can say how a family kept a festival day.}

Complete the last lesson of chapter 328: I can say how a family kept a festival day.
\end{chapteropening}

\section[\texorpdfstring{\sk{व्रतम्} (vratam)}{व्रतम् (vratam)}]{\sk{व्रतम्} (vratam) — a vow; a religious observance}

\label{lesson:SA-C328-vratam}



\begin{quote}
Before the new one: say the Sanskrit for scatters, then the Sanskrit for tosses up.
\end{quote}

\subsection*{You'll want to know: \sk{व्रतम्}}
\textbf{\sk{व्रतम्}} — \emph{vratam} — \textquotedblleft{}a vow; a religious observance\textquotedblright{}.

\begin{quote}
\sk{माता} \sk{अद्य} \sk{व्रतं} \sk{पालयति।}
\end{quote}

\begin{quote}
\emph{mātā adya vrataṃ pālayati.}
\end{quote}

\begin{quote}
Mother keeps a vow today.
\end{quote}

A neuter noun.

\begin{cousinweb}
\textbf{\sk{व्रतम्}} (\emph{vratam}) comes from the root √vṛ, \textquotedblleft{}to choose\textquotedblright{}, by MW's account; MW gives \emph{vrata} as will, command, rule, and also a vow or religious observance (MW).
\end{cousinweb}

\begin{grammarlens}[title={What you've built}]
One more word for telling what happened.
\end{grammarlens}

\subsection*{Your turn}
\begin{itemize}
\raggedright
  \item \emph{Say it:} \emph{vratam}
  \item \emph{Say it:} \emph{vratam}, once more
  \item \emph{Say it:} \emph{mātā adya vrataṃ pālayati.} — Mother keeps a vow today.
  \item \emph{Recall:} say the Sanskrit for scatters, then the Sanskrit for tosses up, then say \emph{vratam} again
\end{itemize}

\begin{tcolorbox}[breakable,colback=teal!4,colframe=teal!35!black,title={Before you move on}]
What does \sk{व्रतम्} mean? (\textquotedblleft{}A vow; a religious observance\textquotedblright{}.) Say it once more.
\end{tcolorbox}
\section[\texorpdfstring{\sk{उपवासः} (upavāsaḥ)}{उपवासः (upavāsaḥ)}]{\sk{उपवासः} (upavāsaḥ) — a fast}

\label{lesson:SA-C328-upavasah}



\begin{quote}
Before the new one: say the Sanskrit for tosses up, then the Sanskrit for a vow.
\end{quote}

\subsection*{You'll want to know: \sk{उपवासः}}
\textbf{\sk{उपवासः}} — \emph{upavāsaḥ} — \textquotedblleft{}a fast\textquotedblright{}.

\begin{quote}
\sk{सा} \sk{सर्वं} \sk{दिनम्} \sk{उपवासं} \sk{करोति।}
\end{quote}

\begin{quote}
\emph{sā sarvaṃ dinam upavāsaṃ karoti.}
\end{quote}

\begin{quote}
She fasts the whole day.
\end{quote}

A masculine noun.

\begin{cousinweb}
\textbf{\sk{उपवासः}} (\emph{upavāsaḥ}) is built from \emph{upa} and the root √vas, \textquotedblleft{}to dwell, stay\textquotedblright{}; MW gives \emph{upavāsa} as a fast, fasting as a religious act (MW).
\end{cousinweb}

\begin{grammarlens}[title={What you've built}]
One more word for telling what happened.
\end{grammarlens}

\subsection*{Your turn}
\begin{itemize}
\raggedright
  \item \emph{Say it:} \emph{upavāsaḥ}
  \item \emph{Say it:} \emph{upavāsaḥ}, once more
  \item \emph{Say it:} \emph{sā sarvaṃ dinam upavāsaṃ karoti.} — She fasts the whole day.
  \item \emph{Recall:} say the Sanskrit for tosses up, then the Sanskrit for a vow, then say \emph{upavāsaḥ} again
\end{itemize}

\begin{tcolorbox}[breakable,colback=teal!4,colframe=teal!35!black,title={Before you move on}]
What does \sk{उपवासः} mean? (\textquotedblleft{}A fast\textquotedblright{}.) Say it once more.
\end{tcolorbox}
\section[\texorpdfstring{\sk{नैवेद्यम्} (naivedyam)}{नैवेद्यम् (naivedyam)}]{\sk{नैवेद्यम्} (naivedyam) — a food offering}

\label{lesson:SA-C328-naivedyam}



\begin{quote}
Before the new one: say the Sanskrit for a vow, then the Sanskrit for a fast.
\end{quote}

\subsection*{You'll want to know: \sk{नैवेद्यम्}}
\textbf{\sk{नैवेद्यम्}} — \emph{naivedyam} — \textquotedblleft{}a food offering\textquotedblright{}.

\begin{quote}
\sk{माता} \sk{मन्दिरे} \sk{नैवेद्यं} \sk{स्थापयति।}
\end{quote}

\begin{quote}
\emph{mātā mandire naivedyaṃ sthāpayati.}
\end{quote}

\begin{quote}
Mother places a food offering in the temple.
\end{quote}

A neuter noun.

\begin{cousinweb}
\textbf{\sk{नैवेद्यम्}} (\emph{naivedyam}) is usually traced to \emph{nivedya}, \textquotedblleft{}to be offered\textquotedblright{}, from the verb of \emph{nivedayati}, with the lengthened first vowel; MW gives \emph{naivedya} as an offering of food presented to a deity (MW).
\end{cousinweb}

\begin{grammarlens}[title={What you've built}]
One more word for telling what happened.
\end{grammarlens}

\subsection*{Your turn}
\begin{itemize}
\raggedright
  \item \emph{Say it:} \emph{naivedyam}
  \item \emph{Say it:} \emph{naivedyam}, once more
  \item \emph{Say it:} \emph{mātā mandire naivedyaṃ sthāpayati.} — Mother places a food offering in the temple.
  \item \emph{Recall:} say the Sanskrit for a vow, then the Sanskrit for a fast, then say \emph{naivedyam} again
\end{itemize}

\begin{tcolorbox}[breakable,colback=teal!4,colframe=teal!35!black,title={Before you move on}]
What does \sk{नैवेद्यम्} mean? (\textquotedblleft{}A food offering\textquotedblright{}.) Say it once more.
\end{tcolorbox}
\section[\texorpdfstring{\sk{अर्पयति} (arpayati)}{अर्पयति (arpayati)}]{\sk{अर्पयति} (arpayati) — offers}

\label{lesson:SA-C328-arpayati}



\begin{quote}
Before the new one: say the Sanskrit for a fast, then the Sanskrit for a food offering.
\end{quote}

\subsection*{You'll want to know: \sk{अर्पयति}}
\textbf{\sk{अर्पयति}} — \emph{arpayati} — \textquotedblleft{}offers\textquotedblright{}.

\begin{quote}
\sk{बालिका} \sk{पुष्पाणि} \sk{अर्पयति।}
\end{quote}

\begin{quote}
\emph{bālikā puṣpāṇi arpayati.}
\end{quote}

\begin{quote}
The girl offers flowers.
\end{quote}

Past (she …): \sk{सा} \sk{पुष्पाणि} \sk{अर्पितवती}, \emph{sā puṣpāṇi arpitavatī}.

\begin{cousinweb}
\textbf{\sk{अर्पयति}} (\emph{arpayati}) comes from the root √ṛ, \textquotedblleft{}to go, move\textquotedblright{}: MW lists \emph{arpaya} as its causative, to make go, hand over, offer (MW).
\end{cousinweb}

\begin{grammarlens}[title={What you've built}]
One more word for telling what happened.
\end{grammarlens}

\subsection*{Your turn}
\begin{itemize}
\raggedright
  \item \emph{Say it:} \emph{arpayati}
  \item \emph{Say it:} \emph{arpayati}, once more
  \item \emph{Say it:} \emph{bālikā puṣpāṇi arpayati.} — The girl offers flowers.
  \item \emph{Recall:} say the Sanskrit for a fast, then the Sanskrit for a food offering, then say \emph{arpayati} again
\end{itemize}

\begin{tcolorbox}[breakable,colback=teal!4,colframe=teal!35!black,title={Before you move on}]
What does \sk{अर्पयति} mean? (\textquotedblleft{}Offers\textquotedblright{}.) Say it once more.
\end{tcolorbox}
\section[\texorpdfstring{\sk{प्रसादः} (prasādaḥ)}{प्रसादः (prasādaḥ)}]{\sk{प्रसादः} (prasādaḥ) — a blessed offering of food; grace}

\label{lesson:SA-C328-prasadah}



\begin{quote}
Before the new one: say the Sanskrit for a food offering, then the Sanskrit for offers.
\end{quote}

\subsection*{You'll want to know: \sk{प्रसादः}}
\textbf{\sk{प्रसादः}} — \emph{prasādaḥ} — \textquotedblleft{}a blessed offering of food; grace\textquotedblright{}.

\begin{quote}
\sk{उत्सवस्य} \sk{अनन्तरं} \sk{सर्वे} \sk{प्रसादं} \sk{खादन्ति।}
\end{quote}

\begin{quote}
\emph{utsavasya anantaraṃ sarve prasādaṃ khādanti.}
\end{quote}

\begin{quote}
After the festival everyone eats the blessed food.
\end{quote}

A masculine noun.

\begin{cousinweb}
\textbf{\sk{प्रसादः}} (\emph{prasādaḥ}) is built from \emph{pra} and the root √sad, \textquotedblleft{}to sit, settle\textquotedblright{}; MW gives \emph{prasāda} as clearness, favour, grace, and also the food presented to a deity and then shared (MW).
\end{cousinweb}

\begin{grammarlens}[title={What you've built}]
One more word for telling what happened.
\end{grammarlens}

\subsection*{Your turn}
\begin{itemize}
\raggedright
  \item \emph{Say it:} \emph{prasādaḥ}
  \item \emph{Say it:} \emph{prasādaḥ}, once more
  \item \emph{Say it:} \emph{utsavasya anantaraṃ sarve prasādaṃ khādanti.} — After the festival everyone eats the blessed food.
  \item \emph{Recall:} say the Sanskrit for a food offering, then the Sanskrit for offers, then say \emph{prasādaḥ} again
\end{itemize}

\begin{tcolorbox}[breakable,colback=teal!4,colframe=teal!35!black,title={Before you move on}]
What does \sk{प्रसादः} mean? (\textquotedblleft{}A blessed offering of food; grace\textquotedblright{}.) Say it once more.
\end{tcolorbox}
